\subsection{Normal Approximation and Confidence Intervals}
% Source integration: supplied Chapter 7 screenshots, Section 7.3.3 and
% selected confidence-interval problems.

The exact design formulas in Theorem \ref{thm:srs-sample-mean-variance} give
the center and scale of the sampling distribution.  For large samples, a normal
approximation turns those formulas into error probabilities and confidence
intervals.

\begin{method} \label{meth:srs-normal-approximation} \hlt{Normal Approximation for SRS}\\
For a simple random sample without replacement, use
\[
\overline{X}\approx N(\mu,\sigma_{\overline{X}}^2),
\qquad
\sigma_{\overline{X}}^2
=\frac{\sigma^2}{n}\left(\frac{N-n}{N-1}\right).
\]
Thus, for $\delta>0$,
\[
\mathbb{P}(\abs{\overline{X}-\mu}\leq \delta)
\approx
2\Phi\left(\frac{\delta}{\sigma_{\overline{X}}}\right)-1.
\]
The same approximation gives
$T-\tau=N(\overline{X}-\mu)\approx N(0,N^2\sigma_{\overline{X}}^2)$.
\end{method}

\begin{definition} \label{def:normal-right-tail-quantile} \hlt{Normal Right-Tail Quantile}\\
For $0<\alpha<1$ and $Z\sim N(0,1)$, $z(\alpha)\defeq\Phi^{-1}(1-\alpha)$, so $\mathbb{P}(Z>z(\alpha))=\alpha$.
\end{definition}

\begin{method} \label{meth:srs-normal-confidence-intervals} \hlt{Approximate SRS Confidence Intervals}\\
For simple random sampling, replace unknown standard errors by the estimates
from Corollary \ref{cor:srs-estimated-standard-errors}.
\begin{enumerate}[label=(\roman*)]
\setlength{\itemsep}{0pt}
\item A two-sided approximate $100(1-\alpha)\%$ confidence interval for $\mu$ is
\[
\overline{X}\pm z(\alpha/2)s_{\overline{X}}.
\]
\item A corresponding interval for the population total $\tau$ is
\[
T\pm z(\alpha/2)s_T.
\]
\item In the dichotomous case, an approximate interval for $p$ is
\[
\widehat p\pm z(\alpha/2)s_{\widehat p}.
\]
\end{enumerate}
The interval is random before sampling.  After the data are observed, the
population parameter is fixed; the coverage statement refers to the long-run
procedure.
\end{method}

\begin{example} \label{ex:srs-condominium-confidence-intervals} \hlt{Condominium Survey Intervals}\\
Suppose a development has $N=8000$ condominium units and a simple random sample
of $n=100$ occupants gives $\overline{X}=1.6$ motor vehicles per unit with
sample standard deviation $s=0.8$.  Then
\[
s_{\overline{X}}
=\frac{0.8}{10}\sqrt{1-\frac{100}{8000}}
\approx 0.08.
\]
Using $z(0.025)\approx1.96$, an approximate $95\%$ interval for the mean number
of vehicles per unit is
\[
1.6\pm1.96(0.08)=(1.44,1.76).
\]
For the total number of vehicles, $T=8000(1.6)=12800$ and
$s_T=8000(0.08)=640$, so the corresponding interval is approximately
$(11546,14054)$.

If the same survey gives $\widehat p=0.12$ for the proportion planning to sell,
then $s_{\widehat p}\approx0.03$, so an approximate $95\%$ interval for $p$ is
$(0.06,0.18)$.  The estimated total number planning to sell is
$8000(0.12)=960$, with estimated standard error about $240$ and approximate
$95\%$ interval $(490,1430)$.
\end{example}

\begin{method} \label{meth:srs-sample-size-from-half-width} \hlt{Sample Size from Desired Half-Width}\\
If a normal-approximation interval has half-width $h$, then ignoring the finite
population correction gives the rough planning rule
\[
h\approx z(\alpha/2)\frac{\sigma}{\sqrt n}
\quad\Longrightarrow\quad
n\approx \left(\frac{z(\alpha/2)\sigma}{h}\right)^2.
\]
For a total based on a proportion estimate, the same calculation uses
$N\sqrt{\widehat p(1-\widehat p)/(n-1)}$ as the estimated standard error.
Thus halving a confidence interval's width requires roughly quadrupling the
sample size when the finite population correction is negligible.
\end{method}
